\begin{table}[H]
\centering
\caption{Campagne d'attaque contre la passerelle \aka (profil LAN)}
\label{tab:attaques-banc}
\small
\begin{tabularx}{\textwidth}{|C{0.06\textwidth}|X|L{0.27\textwidth}|C{0.13\textwidth}|}
\hline
\textbf{Sc.} & \textbf{Attaque} & \textbf{Issue attendue} & \textbf{Conforme}\\ \hline
1 & Rejeu d'une poignée de main enregistrée & \texttt{SIGNATURE\_INVALIDE} & 10/10\\ \hline
2 & Rejeu d'une instruction acceptée et dépassée & \texttt{REJEU\_DETECTE} & 10/10\\ \hline
2 & Relance à l'identique de la dernière instruction (après délai) & \texttt{DOUBLON\_IDEMPOTENT} & 10/10\\ \hline
3 & Altération d'une instruction en transit & \texttt{INTEGRITE\_COMPROMISE} & 10/10\\ \hline
3 & Instruction hors habilitation (plafond) & \texttt{POLITIQUE\_REFUSEE} & 10/10\\ \hline
4 & Réécriture de la liste des suites & \texttt{SIGNATURE\_INVALIDE} & 10/10\\ \hline
-- & Partenaire non inscrit au registre (E1) & \texttt{REGISTRE\_INCONNU} & 10/10\\ \hline
-- & Attaque par réflexion (\S\ref{sec:reflexion}) & \texttt{SIGNATURE\_INVALIDE} & 10/10\\ \hline
\end{tabularx}
\begin{flushleft}\footnotesize
\smallskip
Une attaque est neutralisée si la passerelle la rejette avec le motif attendu, journalisé, et n'accepte aucune instruction adverse. Pour la relance à l'identique, l'issue attendue est la réponse idempotente : le statut initial (acceptée) est renvoyé et le bouchon \emph{core banking} confirme une seule exécution.
 Attaque par réflexion rejouée contre la version~1 du protocole (sans étiquettes de rôle) : instruction forgée acceptée dans 10 cas sur 10 ; contre la version~2 : 10 rejets sur 10 (\texttt{SIGNATURE\_INVALIDE}).
\end{flushleft}
\end{table}
